\documentclass[10pt]{article}
\usepackage[margin=0.8cm]{geometry}
\usepackage{amssymb}
\usepackage{pdfpages}
\usepackage{graphbox}
\usepackage{tikz}
\usepackage{multirow} 
\begin{document}

\begin{samepage}
\begin{center}
  {\Large\bf antilogarithms to base 10}
\end{center}

\pagestyle{empty}
\thispagestyle{empty}

\centerline{\small
\begin{tabular}{|r|c|ccc|ccc|ccc|ccccccccc|}\hline
  &&&&&&&&&&&  \multicolumn{9}{|c|}{ADD} \\
  {\normalsize x}&
  {\normalsize 0}&{\normalsize 1}&{\normalsize 2}&
  {\normalsize 3}&{\normalsize 4}&{\normalsize 5}&
  {\normalsize 6}&{\normalsize 7}&{\normalsize 8}&
  {\normalsize 9}&{\normalsize 1}&{\normalsize 2}&
  {\normalsize 3}&{\normalsize 4}&{\normalsize 5}&
  {\normalsize 6}&{\normalsize 7}&{\normalsize 8}&
  {\normalsize 9}\\ \hline
\vphantom{$A^{A^A}$}
\input{antilog_table_values_page1.txt} % created by antilog_number_maker.R
\hline
\end{tabular}
}

\begin{minipage}{140mm}
\noindent{\small For \rule{0mm}{5mm} example, to calculate the
  antilogarithm of~$0.4915$, use the ``0.49'' row, which gives
  ``3097''.  To incorporate the fourth digit of 5, read off ``4'' from
  its column in the differences, which gives~\mbox{$3097+4=3101$} to
  give a result of $3.101$; $10^{0.4915}=3.101$.}
\end{minipage}
\begin{tikzpicture}[remember picture,overlay]
  \node[anchor=south east,
    xshift=-1cm,
    yshift=1.1cm]
  at (current page.south east){\includegraphics[width=2in]{mono-logo.pdf}};
  \node[anchor=south west,
    xshift=1cm,
    yshift=1.1cm]
  at (current page.south west){{\tt https://github.com/RobinHankin/tables.git}};
\end{tikzpicture}

\end{samepage}
\newpage

\vfill{\ }
\vfill{\ }

\centerline{\small
\begin{tabular}{|r|c|ccc|ccc|ccc|ccccccccc|}\hline
  &&&&&&&&&&&  \multicolumn{9}{|c|}{ADD} \\[1mm]
  {\normalsize x}&
  {\normalsize 0}&{\normalsize 1}&{\normalsize 2}&
  {\normalsize 3}&{\normalsize 4}&{\normalsize 5}&
  {\normalsize 6}&{\normalsize 7}&{\normalsize 8}&
  {\normalsize 9}&{\normalsize 1}&{\normalsize 2}&
  {\normalsize 3}&{\normalsize 4}&{\normalsize 5}&
  {\normalsize 6}&{\normalsize 7}&{\normalsize 8}&
  {\normalsize 9}\\ \hline
  \vphantom{$A^{A^A}$}
  \input{antilog_table_values_page2.txt} % created by antilog_number_maker.R
\hline
\end{tabular}
}

\vfill{\ }
\end{document}

